\section{Checking the hypotheses on actual runs}\label{art:hypothesis-examples}

The examples below check the objective assumptions, comparison fields and
full-BFGS histories together. They distinguish a function satisfying a
convergence theorem from an infinite admissible run on that function, and
from existence for every prescribed initial state. All norms in this appendix
are Euclidean, and an admissible run has the meaning of
Definition~\ref{art:admissible}.

\subsection{A nonsmooth certificate and a full-rank calibration}

\begin{proposition}[The mixed ridge, with an explicit history]\label{art:examples-mixed}
Let $F(u,v)=u^2+|v|$ and $K_b=\{F\le b\}$, where $b>0$.
The function is finite and convex, its unique minimizer is zero, and $K_b$
is compact. For every $0<\eta\le b$, the following data satisfy all the
conditions of Definition~\ref{art:band-certificate} on $K_b$:
\[
 W_\eta(x)=xx^T,\qquad m_\eta=\eta^2,\quad C_\eta=2b,
 \quad L_\eta=0,\quad\delta_\eta=1.
\]
For $b=2/3$ an actual infinite admissible run in $K_b$, with
$(c_1,c_2)=(1/4,1/2)$, is
\[
 x_k=\bigl(0,\tfrac23(-\tfrac12)^k\bigr),\qquad
 g_k=(0,(-1)^k),\qquad
 H_k=\begin{pmatrix}1&0\\0&2^{-k}\end{pmatrix},\qquad \alpha_k=1.
\]
It has $H_0=I$ and converges to zero, although the smallest eigenvalue of
$H_k$ tends to zero.
\end{proposition}
\begin{proof}
The objective assertions follow directly from $F(u,v)=u^2+|v|$.
At a differentiable band point, $g=(2u,\operatorname{sign}v)$ and
$g^Tx=2u^2+|v|\ge F(x)\ge\eta$, which gives visibility.
Write $x=(u,v)$, $z=(r,w)$ and $y=\nabla F(z)-\nabla F(x)$.
If $v,w$ have the same sign, then
\[
 t=2(r-u)^2,\qquad y^Tz=2(r-u)r,\qquad
 (y^Tz)^2\le 2F(z)t.
\]
If their signs differ, put $S=|v|+|w|>0$. Then
\[
 t=2(r-u)^2+2S,\qquad y^Tz=2(r-u)r+2|w|.
\]
Cauchy--Schwarz gives
\[
 (y^Tz)^2\le t\left(2r^2+\frac{2|w|^2}{S}\right)
             \le 2F(z)t\le2bt.
\]
Thus secant control holds for all differentiable endpoint pairs in $K_b$,
in particular for the eligible ones. For an eligible pair, $s\ne0$ and $z=x+s$
give $\Pi_s z=\Pi_s x$, so projected transport holds with equality.
The displayed history has $s_k=-H_kg_k$, $y_k=-2g_k$, and
\[
 d_k=2^{-k},\qquad F(x_k)-F(x_{k+1})=\tfrac13d_k,
 \qquad t_k=2d_k.
\]
Consequently both line-search inequalities hold. The BFGS update leaves
the first diagonal entry unchanged and replaces the second by $2^{-k-1}$;
all iterates are differentiable and nonstationary. This verifies the
complete admissible history and the asserted limit.\qed
\end{proof}

\begin{proposition}[A full-rank certificate in every dimension]\label{art:examples-full-rank}
For $F(x)=\tfrac12\norm{x}^2$ on $\R^n$, $n\ge1$, every positive sublevel
$K_b$ is compact and zero is the unique minimizer. For $0<\eta\le b$,
Definition~\ref{art:band-certificate} holds with
\[
 W_\eta(x)=I,\qquad m_\eta=2\eta,\quad C_\eta=1,
 \quad L_\eta=0,\quad\delta_\eta=1.
\]
The field is positive definite of rank $n$. For any $x_0\ne0$, the history
$x_k=2^{-k}x_0$, $g_k=x_k$, $H_k=I$, $\alpha_k=1/2$ is an infinite
admissible run with $(c_1,c_2)=(1/4,1/2)$, contained in $K_{F(x_0)}$ and
converging to zero.
\end{proposition}
\begin{proof}
Visibility is $g^TIg=\norm{x}^2=2F(x)\ge2\eta$. For any two endpoints,
$y=s$, so $y^TIy=t$. For eligible pairs the constant field gives equality in projected transport.
In the stated history the BFGS update fixes $I$ because $y_k=s_k$.
The actual decrease is $3d_k/4$, and the weak-Wolfe inequality holds with
equality since $g_{k+1}=g_k/2$. Nonzero initial data give nonzero gradients
at every finite index; the other assertions follow from the displayed
geometric decay.\qed
\end{proof}

\subsection{Existence of differentiable weak-Wolfe steps}

\begin{lemma}[An existence criterion for admissible runs]\label{art:examples-run-existence}
Let $F:\R^n\to\R$ be finite, convex and bounded below. Suppose that, for
every differentiability point $x$ and every $p\ne0$, the set
\[
 \{t\in\R:F\text{ is not differentiable at }x+tp\}
\]
is countable. At any differentiable $x$ with $g=\nabla F(x)\ne0$, any
$H\succ0$ and any $0<c_1<c_2<1$, the direction $p=-Hg$ admits a positive
step whose endpoint is differentiable and which satisfies strict Armijo
and strict weak Wolfe. If, in addition, every differentiability point has
nonzero gradient, an infinite admissible run exists for every differentiable
$x_0$, every $H_0\succ0$ and every such pair $(c_1,c_2)$.
\end{lemma}
\begin{proof}
Put $\phi(t)=F(x+tp)$, $D=g^Tp<0$, and
$\sigma(t)=(\phi(t)-\phi(0))/t$ for $t>0$.
Continuity and differentiability at zero give $\sigma(t)\to D$ as
$t\downarrow0$. Boundedness below gives $\sigma(t)>c_1D$ for all
sufficiently large $t$. Since $D<c_2D<c_1D$, continuity supplies a
nonempty open interval on which $c_2D<\sigma(t)<c_1D$.
Choose $\alpha$ there outside the countable exceptional set. The upper
bound is strict Armijo. The supporting inequality at $x+\alpha p$ gives
$\nabla F(x+\alpha p)^Tp\ge\sigma(\alpha)>c_2D$, hence strict weak Wolfe.
Its positive secant curvature preserves positive definiteness under BFGS.
Under the additional nonstationarity hypothesis these facts allow the step
construction to be repeated indefinitely.\qed
\end{proof}

The exceptional set concerns differentiability in the full ambient space,
not just differentiability of $t\mapsto F(x+tp)$. The lemma proves existence
of a sequence of acceptable steps; it does not assert that a prescribed
line-search implementation finds them.

\subsection{A simple family with curved switching}

\begin{proposition}[A two-quadratic family]\label{art:examples-two-quadratics}
For $a,b,c>0$, define
\[
 F_{a,b,c}(u,v)=\max\{a u^2+b v^2+v,\ c u^2+b v^2-v\}.
\]
The function has unique minimizer zero and compact sublevels. It satisfies
the hypotheses of Corollary~\ref{art:planar-convergence} on every positive
sublevel, and its two active branches at zero satisfy the affine-independence
and positive-balance assumptions of Corollary~\ref{art:planar-max}.
For every differentiable initial point, every $H_0\succ0$ and every fixed
$0<c_1<c_2<1$, an infinite admissible run exists. Every infinite admissible
run converges in value and iterates to zero. The switching curve is the
parabola $v=(c-a)u^2/2$, a line when $a=c$.
\end{proposition}
\begin{proof}
Both branches have Hessian at least $2\min\{a,b,c\}I$, and their average
is $(a+c)u^2/2+bv^2$. Thus zero is the unique minimizer, sublevels are compact,
and positive contours have no line segment. The two gradients differ on
the tie curve: their second components differ by two. On that curve their
determinant is
\[
 -2u\bigl(a+c+b(c-a)^2u^2\bigr).
\]
It is nonzero unless $u=0$, when the tie point is zero. At a smooth point
affine criticality means stationarity, hence optimality by convexity.
At a nonzero tie point the two linearly independent gradients have affine
hull excluding zero. Therefore $\mathcal C_{\rm aff}=\{0\}$; away from
that set the switching is regular, and Corollary~\ref{art:planar-convergence}
applies on any positive sublevel. At zero the gradients $(0,1)$ and $(0,-1)$
are affinely independent and balance with weights $1/2$.

Along any line through a differentiability point, the difference between
the two branches is a polynomial of degree at most two, nonzero at the
initial point. Hence there are at most two nondifferentiability points on
that line. A differentiable stationary point would be the unique minimizer,
but zero is not differentiable. Lemma~\ref{art:examples-run-existence}
therefore gives the claimed run existence for every prescribed initial
state and valid constant pair.\qed
\end{proof}

\begin{corollary}[Runs on the printed nonradial example]\label{art:examples-nonradial-runs}
For the objective of Theorem~\ref{art:nonradial-example}, every differentiable
$x_0$ with $F(x_0)<1/100$, every $H_0\succ0$ and every $0<c_1<c_2<1$
admit an infinite admissible run, and every such run converges to zero.
In particular one may take $x_0=(1/1000,0)$, for which
$F(x_0)=1001/10^6<1/100$.
\end{corollary}
\begin{proof}
The objective is bounded below by $u^2+v^2$ and has unique minimizer zero.
At zero the distinct active gradients make it nondifferentiable, so every
differentiable point has nonzero gradient. The restriction of each pairwise
branch difference to a nonconstant line is a nonzero polynomial of degree
at most two. Every line therefore meets the equality sets in finitely many
points. Apply
Lemma~\ref{art:examples-run-existence} and then
Theorem~\ref{art:nonradial-example}. The value at the stated initial point
follows by direct substitution.\qed
\end{proof}

\subsection{Why the run-existence qualification matters}

\begin{proposition}[A covered objective with no differentiable first step]\label{art:examples-no-step}
The function $G(u,v)=\max\{u,-u+v,-u-v\}$ is finite, convex, coercive and
positively homogeneous, with unique minimizer zero. For every $\varepsilon>0$,
the point $x_0=(\varepsilon,0)$ is differentiable. With $H_0=I$, its search
ray has no differentiable positive-step endpoint satisfying weak Wolfe for
any $c_2<1$. Thus no infinite admissible run has these initial data, although
$G$ belongs to the class of Theorem~\ref{art:homogeneous}.
\end{proposition}
\begin{proof}
The function is a maximum of linear forms, and
$G(u,v)\ge|u|$, $2G(u,v)\ge|v|$. These inequalities give coercivity and
the unique minimizer. At $x_0$ the gradient is $(1,0)$, so the ray is
$(\varepsilon-\alpha,0)$. For $0<\alpha<\varepsilon$ its gradient remains
$(1,0)$, and weak Wolfe would require $-1\ge-c_2$, which is false.
For $\alpha\ge\varepsilon$ the endpoint is nondifferentiable: the last
two branches tie with distinct gradients (at equality all three tie).
These cases exhaust the positive ray.\qed
\end{proof}

This is an obstruction to a differentiable infinite history from prescribed
data, not an example of nonoptimal stagnation: the ray itself passes through
the minimizer. The explicit histories above and the existence criterion
specify which stronger non-vacuity conclusions accompany the convergence
examples.
